% GENERATED FILE. Edit canonical lessons, then run npm run generate:books.
% canonical-source-hash: fnv1a64:51ddcf77427f71ab
% canonical-lessons: ZH-C288-jiu, ZH-W288-jiu, ZH-C288-yaoshi, ZH-C288-conditional, ZH-C288-dehua, ZH-C288-depends, ZH-C288-practice

\chapter{If It Rains}
\label{ch:zh-if-it-rains}
\hlchaptermodality{288}

\begin{chapteropening}
\textbf{By the end of this chapter:} \emph{I can say what will happen if something else does.}

Saying what will happen if something else does with \zh{要是} … \zh{就} … and … \zh{的话}, and what it depends on with \zh{看}, built on one new written character: \zh{就}.
\end{chapteropening}

\section[\texorpdfstring{\zh{就} (jiù)}{就 (jiù)}]{\zh{就} (jiù) — then; right away}

\label{lesson:ZH-C288-jiu}



\begin{quote}
A new chapter, and it opens on the small word that joins what comes next to what came first.

\begin{itemize}
\raggedright
  \item \emph{Recall:} say \zh{以后我想当医生} — \emph{yǐhòu wǒ xiǎng dāng yīshēng} — \textbf{R1}, the chapter you have finished
  \item \emph{Recall:} say \zh{我最喜欢中国菜} — \emph{wǒ zuì xǐhuan Zhōngguó cài} — \textbf{R2}
  \item \emph{Recall:} say \emph{jiù}, old — \textbf{R3}, durable
  \item \emph{Recall:} say \zh{马上} — \emph{mǎshàng}, at once — \textbf{R3}
\end{itemize}
\end{quote}

\subsection*{You'll want to know: \zh{就}}
\textbf{\zh{就}} — \emph{jiù} — \textquotedblleft{}then; right away\textquotedblright{}.

\begin{quote}
\textbf{\zh{我马上就来}} \emph{wǒ mǎshàng jiù lái} — I'll be right there.
\end{quote}

\begin{quote}
\textbf{\zh{吃了饭} \zh{我就去}} \emph{chī le fàn, wǒ jiù qù} — Once I've eaten, I'll go.
\end{quote}

It sounds like \emph{jiù}, old, which you know: the same sound, a different word. The next lesson puts it in your hand.

\begin{cousinweb}
\textbf{\zh{就}} (\emph{jiù}) is inherited from Middle Chinese \emph{dzjuwH} (Guangyun, in Baxter's transcription), glossed 'to accomplish', 'to go to meet' and 'to approach'. 'Then; right away' is not given there, so how the word came by it is not settled.
\end{cousinweb}

\begin{grammarlens}[title={What you've built}]
\zh{就} sits after the person, right before the verb: \zh{我马上就来}. It ties what comes next to what came first: \zh{吃了饭}, \zh{我就去}.
\end{grammarlens}

\subsection*{Your turn}
\begin{itemize}
\raggedright
  \item \emph{Say it:} \emph{jiù}
  \item \emph{Say it:} \emph{jiù}, once more
  \item \emph{Say it:} \zh{我马上就来} — \emph{wǒ mǎshàng jiù lái} — I'll be right there.
  \item \emph{Say it:} \zh{吃了饭} \zh{我就去} — \emph{chī le fàn, wǒ jiù qù} — Once I've eaten, I'll go.
  \item \emph{Recall:} say \emph{jiù}, old, then \zh{就}, then: the same sound, two words
\end{itemize}

\begin{tcolorbox}[breakable,colback=teal!4,colframe=teal!35!black,title={Before you move on}]
What does \zh{就} mean? (\textquotedblleft{}Then; right away\textquotedblright{}.) Where does it go? (\textbf{After the person, right before the verb}.)
\end{tcolorbox}
\section[\texorpdfstring{\zh{就} (jiù)}{就 (jiù)}]{\zh{就} — a dot, a bar, \zh{口} and \zh{小}, then a four-stroke right side}

\label{lesson:ZH-W288-jiu}



\begin{quote}
The word first, then the character inside it.

\begin{itemize}
\raggedright
  \item \emph{Recall:} say \emph{jiù} — \textbf{R1}, one lesson back, before it decays
\end{itemize}

\begin{quote}
\zh{文} \zh{口} \zh{小}
\end{quote}

\emph{Read it:} each one aloud, with nothing said first — \emph{wén}, \emph{kǒu}, \emph{xiǎo}
\end{quote}

\subsection*{Script — a dot, a bar, \zh{口} and \zh{小}, then a four-stroke right side}
\hlblockfigure{\detokenize{figures/ZH-W288-jiu-filmstrip.pdf}}{How \zh{就} is written, stroke by stroke}

\begin{quote}
\zh{就}
\end{quote}

\textbf{Four pieces: three you write, and one new one on the right.}

\textbf{A dot and a horizontal}, as \zh{文} opens.

\textbf{\zh{口}, three strokes}, under them.

\textbf{\zh{小}, three strokes}, under \zh{口}: the hooked vertical, then the left dot, then the right dot.

\textbf{The right side, four strokes}: a horizontal; a long left-falling stroke; a vertical that bends right along the foot and hooks up; and last the dot at the upper right.

2 + 3 + 3 + 4 = twelve.

\subsection*{Your turn}
\begin{itemize}
\raggedright
  \item \emph{Write it:} \zh{就}, all twelve, with the model in front of you
  \item \emph{Cover:} write it again from memory, and count the pieces, not the strokes
\end{itemize}

\begin{tcolorbox}[breakable,colback=teal!4,colframe=teal!35!black,title={Before you move on}]
How many strokes in \zh{就}? (\textbf{Twelve}.) Which character you write opens the same way? (\textbf{\zh{文}: a dot and a horizontal}.) Which stroke comes last? (\textbf{The dot at the upper right}.)

Source: \href{https://raw.githubusercontent.com/chanind/hanzi-writer-data/68d10a4b21150cae5e1ebbd223eed289cf32d90c/data/\%E5\%B0\%B1.json}{Hanzi Writer Data \zh{就}, pinned stroke paths and medians}.
\end{tcolorbox}
\section[\texorpdfstring{\zh{要是} (yàoshi)}{要是 (yàoshi)}]{\zh{要是} (yàoshi) — if}

\label{lesson:ZH-C288-yaoshi}



\begin{quote}
Before the new one: the character you wrote last, and a word that sounds the same as this one.

\begin{itemize}
\raggedright
  \item \emph{Read it:} \zh{我马上就来}, aloud — \emph{wǒ mǎshàng jiù lái}
  \item \emph{Recall:} say \emph{yàoshi}, a key — \textbf{R3}, durable
\end{itemize}
\end{quote}

\subsection*{You'll want to know: \zh{要是}}
\textbf{\zh{要是}} — \emph{yàoshi} — \textquotedblleft{}if\textquotedblright{}.

\begin{quote}
\textbf{\zh{要是明天下雨} \zh{我就不去了}} \emph{yàoshi míngtiān xiàyǔ, wǒ jiù bú qù le} — If it rains tomorrow, I won't go.
\end{quote}

It sounds like \emph{yàoshi}, a key: the same sound, a different word. Every character in it is one you write.

\begin{cousinweb}
\textbf{\zh{要是}} (\emph{yàoshi}) is built from \zh{要} and \zh{是}, both inherited from Middle Chinese, \emph{'jiewH} and \emph{dzyeX} (Guangyun, in Baxter's transcription), glossed 'to bind by agreement' and 'right, as against wrong'. 'If' for the pair is the modern word's own sense.
\end{cousinweb}

\begin{grammarlens}[title={What you've built}]
\zh{要是} opens the condition; \zh{就} opens what follows, after its subject: \textbf{\zh{我就不去了}} \emph{wǒ jiù bú qù le}.
\end{grammarlens}

\subsection*{Your turn}
\begin{itemize}
\raggedright
  \item \emph{Say it:} \emph{yàoshi}
  \item \emph{Say it:} \emph{yàoshi}, once more
  \item \emph{Say it:} \zh{要是明天下雨} \zh{我就不去了} — \emph{yàoshi míngtiān xiàyǔ, wǒ jiù bú qù le} — If it rains tomorrow, I won't go.
  \item \emph{Recall:} say \emph{yàoshi}, a key, then \zh{要是}, if
\end{itemize}

\begin{tcolorbox}[breakable,colback=teal!4,colframe=teal!35!black,title={Before you move on}]
What does \zh{要是} mean? (\textquotedblleft{}If\textquotedblright{}.) Which word that sounds the same do you know? (\emph{yàoshi}, a key.)
\end{tcolorbox}
\section[\texorpdfstring{\zh{要是}…\zh{就}… (yàoshi … jiù …)}{要是…就… (yàoshi … jiù …)}]{\zh{要是}…\zh{就}… — if, then}

\label{lesson:ZH-C288-conditional}



\begin{quote}
Two words from this chapter, made into one frame.

\begin{itemize}
\raggedright
  \item \emph{Read it:} \zh{要是}, aloud — \emph{yàoshi}
  \item \emph{Recall:} say \zh{我马上就来} — \emph{wǒ mǎshàng jiù lái} — \textbf{R2}
  \item \emph{Recall:} say \zh{下雨} — \emph{xiàyǔ}, to rain — \textbf{R3}, durable
\end{itemize}
\end{quote}

\begin{grammarlens}[title={Grammar Lens: \zh{要是} … \zh{就} …, if and then}]
\begin{quote}
\textbf{\zh{要是} opens what might happen; \zh{就} opens what follows if it does.}
\end{quote}

\begin{quote}
\textbf{\zh{要是明天下雨} \zh{我就不去了}} \emph{yàoshi míngtiān xiàyǔ, wǒ jiù bú qù le} — If it rains tomorrow, I won't go. \textbf{\zh{要是你病了} \zh{就回家}} \emph{yàoshi nǐ bìng le, jiù huí jiā} — If you're ill, go home.
\end{quote}

\zh{就} goes after the subject of the second clause, right before its verb: \zh{我就不去了}. With no subject there, \zh{就} opens the clause: \zh{就回家}.

\zh{了} at the end of \zh{我就不去了} says the plan has changed.
\end{grammarlens}

\subsection*{Your turn}
Say these aloud:

\begin{itemize}
\raggedright
  \item \zh{要是明天下雨} \zh{我就不去了} — \emph{yàoshi míngtiān xiàyǔ, wǒ jiù bú qù le} — If it rains tomorrow, I won't go.
  \item \zh{要是你病了} \zh{就回家} — \emph{yàoshi nǐ bìng le, jiù huí jiā} — If you're ill, go home.
  \item what you will do tomorrow if it is cold, with \zh{要是} … \zh{就} …
\end{itemize}

\begin{tcolorbox}[breakable,colback=teal!4,colframe=teal!35!black,title={Before you move on}]
Which word opens the condition? (\textbf{\zh{要是}}.) Where does \zh{就} go? (\textbf{After the subject of the second clause, right before its verb}.)
\end{tcolorbox}
\section[\texorpdfstring{\zh{的话} (dehuà)}{的话 (dehuà)}]{\zh{的话} (dehuà) — if, at the end}

\label{lesson:ZH-C288-dehua}



\begin{quote}
Before the new one: last lesson's frame, and a word you have had since the early chapters.

\begin{itemize}
\raggedright
  \item \emph{Read it:} \zh{要是你病了} \zh{就回家}, aloud — \emph{yàoshi nǐ bìng le, jiù huí jiā}
  \item \emph{Recall:} say \zh{话} — \emph{huà}, words — \textbf{R3}, durable
\end{itemize}
\end{quote}

\subsection*{You'll want to know: \zh{的话}}
\textbf{\zh{的话}} — \emph{dehuà} — \textquotedblleft{}if\textquotedblright{}, closing a condition from behind.

\begin{quote}
\textbf{\zh{下雨的话} \zh{我们就不去了}} \emph{xiàyǔ dehuà, wǒmen jiù bú qù le} — If it rains, we won't go.
\end{quote}

\begin{quote}
\textbf{\zh{要是你有时间的话} \zh{就来我家}} \emph{yàoshi nǐ yǒu shíjiān dehuà, jiù lái wǒ jiā} — If you've got time, come round.
\end{quote}

Both characters are ones you write.

\begin{cousinweb}
\textbf{\zh{的话}} (\emph{dehuà}) is built from \zh{的} and \zh{话} 'words'. \zh{话} is inherited from Middle Chinese \emph{hwaejH} (Guangyun, in Baxter's transcription), glossed 'speech'; the particle \zh{的} is the one whose history is not settled. 'If' for the pair is the modern word's own sense, and its first syllable is light.
\end{cousinweb}

\begin{grammarlens}[title={What you've built}]
\zh{要是} opens the condition from the front; \zh{的话} closes it from behind. Use either, or both: \zh{要是你有时间的话}.
\end{grammarlens}

\subsection*{Your turn}
\begin{itemize}
\raggedright
  \item \emph{Say it:} \emph{dehuà}
  \item \emph{Say it:} \emph{dehuà}, once more
  \item \emph{Say it:} \zh{下雨的话} \zh{我们就不去了} — \emph{xiàyǔ dehuà, wǒmen jiù bú qù le} — If it rains, we won't go.
  \item \emph{Say it:} \zh{要是你有时间的话} \zh{就来我家} — \emph{yàoshi nǐ yǒu shíjiān dehuà, jiù lái wǒ jiā} — If you've got time, come round.
  \item \emph{Recall:} say \zh{要是}, then \zh{的话}: the front of a condition, then its back
\end{itemize}

\begin{tcolorbox}[breakable,colback=teal!4,colframe=teal!35!black,title={Before you move on}]
What does \zh{的话} mean? (\textquotedblleft{}If\textquotedblright{}.) Where does it go? (\textbf{At the end of the condition}.)
\end{tcolorbox}
\section[\texorpdfstring{\zh{看天气} (kàn tiānqì)}{看天气 (kàn tiānqì)}]{\zh{看天气} — it depends on the weather}

\label{lesson:ZH-C288-depends}



\begin{quote}
A verb you have used since the first chapters, doing a new job.

\begin{itemize}
\raggedright
  \item \emph{Recall:} say \zh{的话} — \emph{dehuà}, if — \textbf{R1}, one lesson back
  \item \emph{Recall:} say \zh{看} — \emph{kàn}, to look — \textbf{R3}, durable
  \item \emph{Recall:} say \zh{天气} — \emph{tiānqì}, the weather — \textbf{R3}
\end{itemize}
\end{quote}

\begin{grammarlens}[title={Grammar Lens: \zh{看} and a noun, it depends on}]
\begin{quote}
\textbf{\zh{看} and a noun can say 'it depends on': what happens turns on that.}
\end{quote}

\begin{quote}
\textbf{\zh{去不去} \zh{要看天气}} \emph{qù bu qù, yào kàn tiānqì} — Whether we go depends on the weather. \textbf{\zh{那要看你}} \emph{nà yào kàn nǐ} — That's up to you.
\end{quote}

\zh{去不去}, go or not, is the open question; \zh{要看} and the noun says what settles it.
\end{grammarlens}

\subsection*{Your turn}
Say these aloud:

\begin{itemize}
\raggedright
  \item \zh{去不去} \zh{要看天气} — \emph{qù bu qù, yào kàn tiānqì} — Whether we go depends on the weather.
  \item \zh{那要看你} — \emph{nà yào kàn nǐ} — That's up to you.
  \item what your weekend depends on, with \zh{要看}
\end{itemize}

\begin{tcolorbox}[breakable,colback=teal!4,colframe=teal!35!black,title={Before you move on}]
What does \zh{要看天气} say? (\textbf{It depends on the weather}.) And \zh{那要看你}? (\textbf{That's up to you}.)
\end{tcolorbox}
\section[Practice]{\zh{要是下雨} \zh{我就不去了}}

\label{lesson:ZH-C288-practice}



\begin{quote}
The chapter closes here, with a plan that hangs on the weather.

\begin{itemize}
\raggedright
  \item \emph{Read it:} \zh{去不去} \zh{要看天气}, aloud — \emph{qù bu qù, yào kàn tiānqì}
  \item \emph{Recall:} say \zh{的话} — \emph{dehuà} — \textbf{R2}
  \item \emph{Recall:} say \zh{要是} — \emph{yàoshi}, if — \textbf{R2}
\end{itemize}
\end{quote}

\subsection*{The exchange}
\begin{quote}
\textbf{A:} \zh{明天你去山上吗}
\end{quote}

\begin{quote}
\emph{míngtiān nǐ qù shān shàng ma} — Are you going up the mountain tomorrow?
\end{quote}

\begin{quote}
\textbf{B:} \zh{要看天气}
\end{quote}

\begin{quote}
\emph{yào kàn tiānqì} — It depends on the weather.
\end{quote}

\begin{quote}
\textbf{A:} \emph{tiānqì yùbào shuō kěnéng huì xiàyǔ} — The forecast says it might rain.
\end{quote}

\begin{quote}
\textbf{B:} \zh{要是下雨} \zh{我就不去了}
\end{quote}

\begin{quote}
\emph{yàoshi xiàyǔ, wǒ jiù bú qù le} — If it rains, I won't go.
\end{quote}

\begin{quote}
\textbf{A:} \zh{不下雨的话呢}
\end{quote}

\begin{quote}
\emph{bú xiàyǔ dehuà ne} — And if it doesn't?
\end{quote}

\begin{quote}
\textbf{B:} \zh{不下雨的话} \zh{我就去}
\end{quote}

\begin{quote}
\emph{bú xiàyǔ dehuà, wǒ jiù qù} — If it doesn't, I'll go.
\end{quote}

\begin{quote}
\textbf{A:} \zh{要是你有时间的话} \zh{就来我家}
\end{quote}

\begin{quote}
\emph{yàoshi nǐ yǒu shíjiān dehuà, jiù lái wǒ jiā} — If you've got time, come round.
\end{quote}

\begin{quote}
\textbf{B:} \zh{好} \zh{我马上就来}
\end{quote}

\begin{quote}
\emph{hǎo, wǒ mǎshàng jiù lái} — All right, I'll be right over.
\end{quote}

\zh{要是} opens one condition, \zh{的话} closes the next, and \zh{就} opens each result. \zh{要看天气} answers the question before it is settled.

\begin{tcolorbox}[breakable,skin=enhanced,colback=orange!6,colframe=orange!45!black,boxrule=0.5pt,arc=1mm,left=6pt,right=6pt,top=4pt,bottom=4pt,fonttitle=\bfseries,title={Read this}]
\emph{Read it:} the exchange through once without stopping, then answer the question below
\end{tcolorbox}

\begin{grammarlens}[title={What you've built — if, and what it depends on}]
Three ways to put a condition, and one to say it is not settled:

\begin{itemize}
\raggedright
  \item \textbf{\zh{要是下雨} \zh{我就不去了}} \emph{yàoshi xiàyǔ, wǒ jiù bú qù le} — \zh{要是} at the front.
  \item \textbf{\zh{下雨的话} \zh{我们就不去了}} \emph{xiàyǔ dehuà, wǒmen jiù bú qù le} — \zh{的话} at the back.
  \item \textbf{\zh{要是你有时间的话}} \emph{yàoshi nǐ yǒu shíjiān dehuà} — both.
  \item \textbf{\zh{要看天气}} \emph{yào kàn tiānqì} — it depends.
\end{itemize}

\textbf{Listen.} \emph{Yàoshi míngtiān hěn lěng, nǐ qù ma?}

\textbf{Speak.} Say two things you will do this weekend, each depending on something.

\emph{Write it:} \zh{就} — the opening of \zh{文}, \zh{口} and \zh{小}, then the right side
\end{grammarlens}

\subsection*{Your turn}
\begin{itemize}
\raggedright
  \item \emph{Say it:} \zh{要是你有时间的话} \zh{就来我家} — \emph{yàoshi nǐ yǒu shíjiān dehuà, jiù lái wǒ jiā} — If you've got time, come round.
  \item \emph{Say it:} \zh{去不去} \zh{要看天气} — \emph{qù bu qù, yào kàn tiānqì} — Whether we go depends on the weather.
  \item \emph{Say it:} the whole exchange, both parts
  \item \emph{Write it:} \zh{要是下雨} \zh{我就不去了}
\end{itemize}

\begin{tcolorbox}[breakable,colback=teal!4,colframe=teal!35!black,title={Before you move on}]
Say if it rains, you won't go. (\textbf{\zh{要是下雨} \zh{我就不去了}}.) Say it depends on the weather. (\textbf{\zh{要看天气}}.) Say you'll be right there. (\textbf{\zh{我马上就来}}.)
\end{tcolorbox}
